% Generated draft: requires mathematical and source review.
\documentclass[11pt]{article}
\usepackage[margin=25mm]{geometry}
\usepackage{fontspec}
\setmainfont{Latin Modern Roman}
\usepackage{amsmath,amssymb,mathtools,unicode-math}
\setmathfont{Latin Modern Math}
\usepackage{enumitem,needspace,xurl,hyperref}
\setlength{\parindent}{0pt}
\setlength{\parskip}{5pt}
\setlength{\emergencystretch}{3em}
\newcommand{\N}{\mathbb{N}}
\newcommand{\Z}{\mathbb{Z}}
\newcommand{\Q}{\mathbb{Q}}
\newcommand{\R}{\mathbb{R}}
\newcommand{\C}{\mathbb{C}}
\newcommand{\F}{\mathbb{F}}
\newcommand{\A}{\mathbb{A}}
\newcommand{\PP}{\mathbb{P}}
\newcommand{\E}{\mathbb{E}}
\newcommand{\eps}{\varepsilon}
\newcommand{\dd}{\,\mathrm{d}}
\begin{document}
\section*{UnsolvedMath ID: 30003417}
\textbf{Source problem code:} catalogue entry
\subsection{Definitions and mathematical statement}
Let \(\kappa\) be an uncountable regular cardinal, meaning \(\kappa>\aleph_0\) and \(\operatorname{cf}(\kappa)=\kappa\), and assume
\[
2^{<\kappa}:=\sup_{\lambda<\kappa}2^\lambda=\kappa.
\]
For \(f,g\in\kappa^\kappa\), define eventual domination by
\[
f\leq^* g\quad\Longleftrightarrow\quad
\exists\alpha<\kappa\ \forall\beta\in[\alpha,\kappa)\; f(\beta)\leq g(\beta).
\]
Define the bounding and dominating numbers by
\[
\mathfrak b_\kappa=
\min\bigl\{|B|:B\subseteq\kappa^\kappa\ \text{and}\ \neg\exists g\in\kappa^\kappa\ \forall f\in B\;(f\leq^*g)\bigr\},
\]
\[
\mathfrak d_\kappa=
\min\bigl\{|D|:D\subseteq\kappa^\kappa\ \text{and}\ \forall f\in\kappa^\kappa\ \exists g\in D\;(f\leq^*g)\bigr\}.
\]
The supplied record does not define the underlying space or the precise convention for \(\mathcal M_\kappa\). The following definitions therefore provide a conditional formulation. For a specified choice \(A=\{0,1\}\) or \(A=\kappa\), put \(X=A^\kappa\) and give \(X\) the topology with basic open sets
\[
[s]=\{x\in A^\kappa:x\restriction\alpha=s\},
\qquad \alpha<\kappa,\quad s\in A^\alpha.
\]
A set \(N\subseteq X\) is nowhere dense if \(\operatorname{int}_X(\overline N^{\,X})=\varnothing\). Under the convention that a \(\kappa\)-meager set is covered by at most \(\kappa\) nowhere dense sets, define
\[
\mathcal M_\kappa^X=
\left\{B\subseteq X:\exists\langle N_\xi:\xi<\kappa\rangle\;
\left[\forall\xi<\kappa\;(N_\xi\text{ is nowhere dense in }X)
\ \text{and}\ B\subseteq\bigcup_{\xi<\kappa}N_\xi\right]\right\}.
\]
For an ideal \(\mathcal I\) on \(X\), its additivity and cofinality are
\[
\operatorname{add}(\mathcal I)=
\min\left\{|\mathcal F|:\mathcal F\subseteq\mathcal I,
\ \bigcup\mathcal F\notin\mathcal I\right\},
\]
\[
\operatorname{cof}(\mathcal I)=
\min\left\{|\mathcal F|:\mathcal F\subseteq\mathcal I,
\ \forall B\in\mathcal I\ \exists C\in\mathcal F\;(B\subseteq C)\right\}.
\]
\textbf{Question.} For every uncountable regular cardinal \(\kappa\) satisfying \(2^{<\kappa}=\kappa\), must both equalities
\[
\operatorname{add}(\mathcal M_\kappa)=\mathfrak b_\kappa
\qquad\text{and}\qquad
\mathfrak d_\kappa=\operatorname{cof}(\mathcal M_\kappa)
\]
hold?
Here identifying \(\mathcal M_\kappa\) with the explicitly defined \(\mathcal M_\kappa^X\) requires confirmation of the intended space and covering convention. The conditional definitions above do not assert that the two choices of \(X\) give equivalent questions.
\subsection{Short English statement}
For every uncountable regular cardinal \(\kappa\) with \(2^{<\kappa}=\kappa\), does the additivity of the generalized meager ideal equal the bounding number, and does its cofinality equal the dominating number? The source leaves the underlying space and meagerness convention unspecified, so the explicit formulation requires confirmation of those conventions.
\subsection{Sources}
\begin{itemize}
\item Supplied local UnsolvedMath registry record, permanent ID \(30003417\), source problem code \texttt{catalogue entry}, titled ``Meager Ideal Equalities at Uncountable Regular Cardinals,'' in the Oberwolfach Reports Open Problems collection. No individual proposer is supplied; the proposed year is recorded as 2017.
\item Supplied report citation: \emph{Set Theory (2018)}, Oberwolfach Reports, pp.\ 527--589, DOI \texttt{10.4171/owr/2017/11}. \url{https://doi.org/10.4171/owr/2017/11}. The supplied background describes the collection as drawn from workshop reports and their open-problem discussions.
\item Supplied literature reference: \emph{Higher Cicho\'n diagram in the degenerate case}. \url{https://arxiv.org/abs/1907.03111}.
\item Supplied literature reference: \emph{Generalized cardinal invariants with compactness}. \url{https://doi.org/10.1007/s00153-025-00977-2}.
\item Supplied literature reference: \emph{OWR Set Theory report}. \url{https://ems.press/content/serial-article-files/46674}.
\item Literature assessment dated \texttt{2026-08-21}: the supplied curation labels the problem open and attributes to Brendle exact minimum and maximum formulas involving the meager ideal, rather than the two equalities asked here. It also reports equality in special compactness models, without a general result for every regular uncountable cardinal under the stated arithmetic. This is supplied dated curation derived from the OWR collection, not independently verified here; the record itself requests independent verification of time-sensitive status and cited claims.
\item The registry labels the statement \texttt{corrected\_verified}, explaining that the mathematical content was preserved while notation, typography, and sentence structure were repaired for publication. This supplied editorial label does not constitute independent mathematical verification of this draft.
\end{itemize}
\end{document}
